% Originally: content/week-07.tex, \section{Submanifolds} (Corsin Week 7)

% Source: Corsin Nick/Class Notes/Week 7.pdf, p. 2
\section{Submanifolds}
\label{sec:submanifolds}

Recall that a \newterm{$C^k$-diffeomorphism} (\cref{thm:inverse_function_theorem}) is a $C^k$
bijective function whose inverse is also $C^k$.

\begin{definition}[Submanifold]
\label{def:submanifold}
$M^m \subseteq \mathbb{R}^n$ is a $C^k$, $m$-dimensional \newterm{submanifold}
\germanfor{Untermannigfaltigkeit} of $\mathbb{R}^n$ if for any point $p \in M$ there are an
open neighbourhood $U \subseteq \mathbb{R}^n$ of $p$, an open neighbourhood $V \subseteq
\mathbb{R}^n$ of $0$, and a $C^k$-diffeomorphism $\Phi : U \to V$ with $\Phi(p) = 0$ (a
\newterm{submanifold chart}) such that
\[\Phi(U \cap M) = (\mathbb{R}^m\times\{0\}) \cap V = \{x \in V : x_{m+1}=\dots=x_n=0\}.\]
\end{definition}

% Sonnet 5 (Medium) --- FIG-W07-01
\begin{center}
\begin{tikzpicture}[>=Latex, scale=1.0]
  % Left panel: the curved submanifold M with a chart neighbourhood U around p
  \begin{scope}[shift={(0,0)}]
    \draw[blue!70!black, thick] plot[smooth] coordinates
      {(-1.6,-0.4) (-0.9,0.5) (0.0,0.1) (0.9,0.7) (1.7,-0.1)};
    \node[blue!70!black, font=\small] at (1.9,-0.4) {$M$};
    \draw[orange!90!black, thick, dashed] (0.0,0.35) ellipse (0.85 and 0.75);
    \node[orange!90!black, font=\small] at (0.0,1.25) {$U$};
    \fill (0.0,0.1) circle (1.5pt) node[below, font=\small] {$p$};
    \node[below, font=\small] at (0,-1.3) {$\mathbb{R}^n$};
  \end{scope}

  \draw[->, thick] (2.5,0) -- (3.7,0) node[midway, above, font=\small] {$\Phi$};

  % Right panel: the flattened image, R^m x {0} as a horizontal line inside V
  \begin{scope}[shift={(4.4,0)}]
    \draw[orange!90!black, thick] (-1.1,-0.9) rectangle (1.9,1.1);
    \node[orange!90!black, font=\small] at (0.4,1.35) {$V$};
    \draw[green!60!black, thick, ->] (-1.05,0) -- (1.85,0)
      node[right, font=\scriptsize] {$\mathbb{R}^m\times\{0\}$};
    \fill (0,0) circle (1.5pt) node[below, font=\small] {$0 = \Phi(p)$};
    \node[below, font=\small] at (0.4,-1.3) {$\mathbb{R}^n$};
  \end{scope}
\end{tikzpicture}
\end{center}

% Corsin's own observation; the note it used to sit in was an ainote until 2026-08-09, but by
% Test 1 in style.md it is mathematics rather than commentary about the transcription.
\begin{remark}[Why the definition is almost never used directly]
Producing such a $\Phi$ by hand is difficult even for simple $M$, so the definition is rarely
what one works with. Two equivalent descriptions do the job instead, the \textbf{implicit} one
(the regular value theorem below) and the \textbf{explicit or graphical} one (the local
parametrization theorem below). \Cref{rem:three_submanifold_descriptions} compares all three
once both are available.
\end{remark}

\subsection{The regular value theorem}

\begin{theorem}[Regular value theorem]
\label{thm:regular_value_theorem}
Suppose $U \subseteq \mathbb{R}^n$ is open, $F \in C^k(U,\mathbb{R}^\ell)$, and $y \in F(U)$. If
\[DF_x : \mathbb{R}^n \to \mathbb{R}^\ell\]
is surjective (i.e., $\Jac F(x)$ has full rank) for all $x \in F^{-1}\{y\} = \{\bar x \in U :
F(\bar x) = y\}$, then $F^{-1}\{y\}$ is a $C^k$ submanifold of dimension $n-\ell$.
\end{theorem}

% Sonnet 5 (Medium)
\begin{importantremark}[Full rank means two opposite things]
\label{rem:full_rank_two_meanings}
\qt{Full rank} is decided by the \textbf{shape} of the matrix, and the two representations of a
submanifold sit on opposite sides of it:
\begin{itemize}
  \item An \textbf{implicit} description $F : \mathbb{R}^n \to \mathbb{R}^\ell$ cuts $M$ out by
    $\ell$ equations, so $\Jac F$ is \textbf{wide} ($\ell \times n$ with $\ell \leq n$). Full
    rank $=\ell$ means $DF_x$ is \textbf{surjective}, so the $\ell$ constraints are independent,
    each one genuinely cutting a dimension. Here $\dim M = n - \ell$.
  \item A \textbf{parametrization} $f : \mathbb{R}^m \to \mathbb{R}^n$ sweeps $M$ out by $m$
    parameters, so $\Jac f$ is \textbf{tall} ($n \times m$ with $m \leq n$). Full rank $=m$
    means $Df_x$ is \textbf{injective}, so the $m$ parameters move in independent directions,
    none of them wasted. Here $\dim M = m$.
\end{itemize}
Same phrase, opposite condition. A wide matrix is never injective and a tall one is never
surjective, so there is no ambiguity in practice. But it is worth checking which side you are on
before writing down what \qt{full rank} is supposed to give you.
\end{importantremark}

% Source: Corsin Nick/Class Notes/Week 7.pdf, p. 2
\begin{example}[Circle of radius $R$]
Let $F : \mathbb{R}^2 \to \mathbb{R}$, $(x,y) \mapsto x^2+y^2$. Then
\[\Jac F(x,y) = (2x,\ 2y)\]
has full rank at every $(x,y) \in \mathbb{R}^2\setminus\{0\}$. Here $\ell = 1$, so by
\cref{rem:full_rank_two_meanings} full rank means that $DF_{(x,y)} : \mathbb{R}^2 \to \mathbb{R}$
is \textbf{surjective}, which for a single row is simply the requirement that $(2x,2y)$ be
nonzero. For $R>0$, $(0,0) \notin F^{-1}\{R^2\}$, so
\[F^{-1}\{R^2\} = \{(x,y) \in \mathbb{R}^2 : x^2+y^2=R^2\}\]
is a submanifold of dimension $2-1=1$ by \cref{thm:regular_value_theorem}. Similarly,
$\{x \in \mathbb{R}^n : |x|^2=R^2\}$ is a submanifold of dimension $n-1$ for $n \geq 1$ (the
\newterm{$n$-sphere} of radius $R$).
\end{example}

% Sonnet 5 (Medium) --- FIG-W07-02
\begin{center}
\begin{tikzpicture}[>=Latex, scale=1.1]
  \draw[->] (-1.7,0) -- (1.7,0) node[right, font=\small] {$x$};
  \draw[->] (0,-1.7) -- (0,1.7) node[above, font=\small] {$y$};
  \draw[blue!70!black, thick] (0,0) circle (1.3);
  \node[blue!70!black, font=\small] at (1.05,1.05) {$F^{-1}\{R^2\}$};
  \coordinate (Q) at (0.92,0.92);
  \fill (Q) circle (1.5pt);
  \draw[orange!90!black, thick, ->] (Q) -- ++(0.55,0.55)
    node[above right, font=\scriptsize] {$\nabla F(x,y) = (2x,2y)$};
\end{tikzpicture}
\end{center}

% Generator: Claude Opus 5 (Medium)
\begin{aiexample}[A set that is \emph{not} a submanifold, and how to prove it]
\label{ex:cross_not_submanifold}
The regular value theorem is a one-way street: if $DF_x$ is surjective everywhere on the level
set, we get a submanifold. When it is not, we get \emph{no information}. The level set may still
be a submanifold, or it may not. It is worth seeing a case where it is not.

Let $F(x,y) := xy$ and consider $C := F^{-1}\{0\}$, the union of the two coordinate axes.
Here $\Jac F(x,y) = (y,\ x)$, which is nonzero, and hence surjective, at every point of $C$
\emph{except} the origin. So, \cref{thm:regular_value_theorem} does apply on
$C \setminus \{(0,0)\}$, telling us that the four open half-axes are $1$-dimensional
submanifolds, and says nothing at the origin.

And at the origin $C$ genuinely fails to be a submanifold. The argument is a connectedness trick
worth remembering, since it is the standard way to prove a negative here.

Suppose $C$ were a $1$-dimensional submanifold. By \cref{def:submanifold} there would be a
neighbourhood $U$ of the origin and a diffeomorphism $\Phi : U \to V$ carrying $C \cap U$ onto a
piece of a \emph{line}, $(\mathbb{R}\times\{0\})\cap V$. Now delete the centre point from each
side and count connected components (\cref{def:connectedness}):
\begin{itemize}
  \item Removing a point from a small piece of a line leaves \textbf{two} pieces.
  \item Removing the origin from $C \cap U$ leaves \textbf{four} pieces, one along each
    half-axis.
\end{itemize}
But $\Phi$ is a homeomorphism, and homeomorphisms preserve the number of connected components
(\cref{thm:continuity_preserves_connected}, applied to $\Phi$ and to $\Phi^{-1}$). Four cannot
equal two, so no such $\Phi$ exists and $C$ is not a submanifold.

\begin{remark}
The moral is that submanifolds are forbidden from having \emph{crossings, corners or cusps}: near
each of its points a submanifold must look like a flat piece of $\mathbb{R}^m$, and a crossing
does not. Compare \cref{ex:constraint_qualification_fails}, where the constraint set
$\{y^2=x^3\}$ had a cusp at exactly the point where $\nabla g$ vanished. That is the same
phenomenon, seen through the Lagrange-multiplier lens back in \cref{ch:lagrange}.
\end{remark}
\end{aiexample}

% Source: Corsin Nick/Class Notes/Week 7.pdf, p. 4
\subsubsection{Connection to Lagrange multipliers}

\begin{remark}[Constraint sets are submanifolds]
Recall the Lagrangian from \cref{prop:constrained_optimization}: when optimizing $f|_M$ for
$f \in C^k(\mathbb{R}^n,\mathbb{R})$, where $g_1,\dots,g_k : U \to \mathbb{R}$,
\[M := \{x \in \mathbb{R}^n : g_1(x)=\dots=g_k(x)=0\}, \qquad (\nabla g_1,\dots,\nabla g_k) \text{ linearly independent for all } x \in M\]
(hence $k \leq n$). This means precisely that, for $g : \mathbb{R}^n \to \mathbb{R}^k$,
$x \mapsto (g_1(x),\dots,g_k(x))$, the Jacobian
\[\Jac g(x) = \begin{pmatrix} \nabla g_1(x) \longrightarrow \\ \vdots \\ \nabla g_k(x) \longrightarrow \end{pmatrix}\]
has full rank. Therefore, $M = g^{-1}\{0\}$ is an $(n-k)$-dimensional submanifold by the regular
value theorem: the constraint sets of the Lagrange-multiplier problems in \cref{ch:lagrange}
were submanifolds all
along.
\end{remark}

% Source: Corsin Nick/Class Notes/Week 7.pdf, p. 5
\subsection{The local parametrization theorem}

\begin{theorem}[Local parametrization]
\label{thm:local_parametrization}
Let $M \subseteq \mathbb{R}^n$. Then $M$ is a $C^k$ submanifold with $\dim M = m$ if and only if
for every $p \in M$ there exist an open neighbourhood $V$ of $p$, an open set $U \subseteq
\mathbb{R}^m$, and a $C^k$ map $f : U \to V \cap M$ such that
\begin{itemize}
  \item $Df_x : \mathbb{R}^m \to \mathbb{R}^n$ is injective for all $x \in U$;
  \item $f$ is a homeomorphism (bijective, continuous, with continuous inverse $f^{-1}$).
\end{itemize}
\end{theorem}

% Generator: Claude Opus 5 (Medium)
\begin{remark}[Three descriptions, and which one to reach for]
\label{rem:three_submanifold_descriptions}
There are now three ways to certify that $M$ is an $m$-dimensional submanifold, and they are
equivalent. They are not equally convenient, though, and choosing badly is what makes these
exercises feel hard.

\begin{itemize}
  \item \textbf{Chart} (\cref{def:submanifold}) is the definition itself. Straighten $M$ out
    with a diffeomorphism, so that the first $m$ coordinates of the chart carry the submanifold
    and the rest vanish. It is almost never used directly on a concrete $M$, since producing
    $\Phi$ by hand is painful. Its role is theoretical, and it is what the other two are proved
    from.
  \item \textbf{Implicit, via the regular value theorem} (\cref{thm:regular_value_theorem}) is
    the one to reach for when $M$ arrives as \emph{equations}, such as $\{x^2+y^2=R^2\}$ or
    $\{g_1=\dots=g_k=0\}$. It rests on the same idea as the chart definition, with the
    diffeomorphism produced implicitly rather than by hand. Check a rank condition on $\Jac F$
    and you are done, with no parametrization ever written. This is normally the cheapest route.
  \item \textbf{Parametrization} (\cref{thm:local_parametrization}) is the one for an $M$ that
    arrives as a \emph{picture you can sweep out}, such as a graph, a surface of revolution, or
    a curve given by a formula in a parameter. It is often the most intuitive, since it builds
    the submanifold out of parameters explicitly. It is also the one to use when the tangent
    space is wanted next, because $\partial_1f,\dots,\partial_mf$ hand you a basis directly
    (\cref{def:tangent_space}).
\end{itemize}

The rough rule: \textbf{equations $\to$ regular value theorem; a formula with parameters $\to$
local parametrization.} A graph $\{(x,f(x))\}$ satisfies both, which is why graphs are the
easiest case of all.

Two cautions carried over from \cref{rem:full_rank_two_meanings}. \qt{Full rank} means
\emph{surjective} in the first case and \emph{injective} in the second, and the shape of the
matrix decides which. And neither theorem yields a negative. If the rank condition fails you
have learned nothing, and must argue separately (see \cref{ex:cross_not_submanifold}).
\end{remark}

\begin{example}[The $n$-parabola]
Let $P := \{(x,t) \in \mathbb{R}^n\times\mathbb{R} : t=|x|^2\}$. One could use the regular value
theorem, but instead take
\[f : \mathbb{R}^n \to \mathbb{R}^n\times\mathbb{R}, \qquad x \mapsto (x,|x|^2).\]
\begin{enumerate}[label=\textbf{(\alph*)}]
  \item The Jacobian matrix has full rank $n$ at every $x$:
\[ \Jac f(x) = \begin{pmatrix} \id_{n\times n} \\ 2x\transp \end{pmatrix}. \]
    This one is tall, so by \cref{rem:full_rank_two_meanings} full rank means that $Df_x$ is
    \textbf{injective}, which is clear from the $\id_{n\times n}$ block alone.
  \item $f$ is bijective onto $P$ and continuous, with inverse $\pi : f(U) \to \mathbb{R}^n$,
    $(x,t) \mapsto x$ (the canonical projection), which is also continuous.
  \item $f(\mathbb{R}^n) = P$.
\end{enumerate}
So, $P$ is an $n$-dimensional submanifold by \cref{thm:local_parametrization}.

\begin{remark}
The graphical representation of a submanifold (as the graph of a function, here $x \mapsto |x|^2$)
is therefore a common, and equivalent, special case of a local parametrization: it is always
injective on its domain by construction, and its inverse is simply the projection forgetting the
last coordinate.
\end{remark}
\end{example}

% Sonnet 5 (Medium) --- FIG-W07-03
\begin{center}
\begin{tikzpicture}[>=Latex, scale=1.1]
  \draw[->] (-1.9,0) -- (1.9,0) node[right, font=\small] {$x$};
  \draw[->] (0,-0.3) -- (0,2.7) node[above, font=\small] {$t$};
  \draw[blue!70!black, thick, domain=-1.6:1.6, samples=60, smooth]
    plot ({\x}, {\x*\x}) node[above right, font=\small] {$P = \{(x,t) : t=|x|^2\}$};
\end{tikzpicture}
\end{center}

% Source: Corsin Nick/Class Notes/Week 7.pdf, p. 6
\subsubsection{An alternative optimization technique}

\begin{remark}[Parametrized optimization, revisiting \cref{ex:extrema_circle_disk}]
Local parametrizations give a second route to constrained-extremum problems, alongside Lagrange
multipliers: parametrize the constraint set and optimize the resulting unconstrained function of
fewer variables. Corsin illustrates this by revisiting $f(x,y) := 2x^2+y^2-x$ on
$S^1 = \{(x,y)\in\mathbb{R}^2 : x^2+y^2=1\}$, which is exactly \cref{ex:extrema_circle_disk}\textbf{(a)}, solved
there via Lagrange multipliers.

Parametrize $S^1$ by $(x(\theta),y(\theta)) := (\cos\theta,\sin\theta)$. Then
\[f(\theta) := f(x(\theta),y(\theta)) = 2\cos^2\theta+\sin^2\theta-\cos\theta = \cos^2\theta-\cos\theta+1,\]
using $\sin^2\theta = 1-\cos^2\theta$. Differentiating,
\[f'(\theta) = -2\cos\theta\sin\theta+\sin\theta \overset{!}{=} 0.\]
\begin{enumerate}[label=\textbf{(\alph*)}]
  \item $\sin\theta = 0 \implies \theta_1 \in \{0,\pi\}$.
  \item $\sin\theta \neq 0 \implies -2\cos\theta+1=0 \implies \cos\theta = \tfrac12 \implies
    \theta_2 \in \{\tfrac{\pi}{3},\tfrac{5\pi}{3}\}$.
\end{enumerate}
So, the extrema are at $(\cos\theta_1,\sin\theta_1) = (\pm1,0)$ and
$(\cos\theta_2,\sin\theta_2) = \big(\tfrac12,\pm\tfrac{\sqrt3}{2}\big)$. These are exactly the
four points found via Lagrange multipliers in the solution to \cref{ex:extrema_circle_disk}\textbf{(a)}, so the
two methods agree.
\end{remark}

% Source: Jérôme Paschoud/Notizen/Woche 7.2 Untermannigfaltigkeiten.pdf, p. 4
% Generator: Gemini 3.6 Flash (High)
\begin{exercise}[Dimension of a submanifold]
\label{ex:dimension_submanifold_system}
Show that $M := \{(x,y,z) \in \mathbb{R}^3 \mid 2x+z=1, x^2+y^2-z^2=1\}$ is a submanifold of $\mathbb{R}^3$. What is its dimension? (cf.\ Michaels Bsp 9.1.2 d)
\end{exercise}

\begin{exercisesolution}[\cref{ex:dimension_submanifold_system}]
Let $F(x,y,z) = \begin{pmatrix} 2x+z-1 \\ x^2+y^2-z^2-1 \end{pmatrix}$. Then $M = F^{-1}\{0\}$.
We compute the Jacobian matrix of $F$:
\[ \Jac F(x,y,z) = \begin{pmatrix} 2 & 0 & 1 \\ 2x & 2y & -2z \end{pmatrix}. \]
We need to check that $\rank(\Jac F(x,y,z)) = 2$ for all $(x,y,z) \in M$.
The rank is less than 2 if and only if all $2 \times 2$ subdeterminants are zero:
\begin{align*}
\det \begin{pmatrix} 2 & 0 \\ 2x & 2y \end{pmatrix} &= 4y = 0 \implies y = 0 \\
\det \begin{pmatrix} 2 & 1 \\ 2x & -2z \end{pmatrix} &= -4z - 2x = 0 \implies x = -2z \\
\det \begin{pmatrix} 0 & 1 \\ 2y & -2z \end{pmatrix} &= -2y = 0 \implies y = 0.
\end{align*}
So, the rank is not maximal only when $y=0$ and $x=-2z$.
Let's check if such a point can belong to $M$. Substituting $x=-2z$ into the first equation of $M$ gives $2(-2z) + z = -3z = 1 \implies z = -1/3$.
Then $x = 2/3$. Thus the point is $(2/3, 0, -1/3)$.
Now check the second equation of $M$: $(2/3)^2 + 0^2 - (-1/3)^2 = 4/9 - 1/9 = 3/9 = 1/3 \neq 1$.
Therefore, this point is not in $M$.
Since $\Jac F(x,y,z)$ has maximal rank 2 for all $(x,y,z) \in M$, $M$ is a submanifold of $\mathbb{R}^3$.
Its dimension is $n - k = 3 - 2 = 1$.
\end{exercisesolution}

% Source: Jérôme Paschoud/Notizen/Woche 7.2 Untermannigfaltigkeiten.pdf, p. 4
% Generator: Gemini 3.6 Flash (High)
\begin{exercise}[Submanifold of matrices]
\label{ex:submanifold_matrices_rank1}
Show that $M := \{ A \in \operatorname{Mat}_{2\times 2}(\mathbb{R}) \mid \rank(A) = 1 \}$ is a submanifold of $\operatorname{Mat}_{2\times 2}(\mathbb{R})$. What is the dimension of $M$?
\end{exercise}

\begin{exercisesolution}[\cref{ex:submanifold_matrices_rank1}]
Let $A = \begin{pmatrix} a & b \\ c & d \end{pmatrix}$. The condition $\rank(A) \leq 1$ is equivalent to $\det(A) = ad - bc = 0$.
The condition $\rank(A) = 1$ means $\det(A) = 0$ but $A \neq 0$.
So, $M$ is the zero set of $F : \mathbb{R}^4 \setminus \{0\} \to \mathbb{R}$, $F(a,b,c,d) = ad - bc$.
The space $\operatorname{Mat}_{2\times 2}(\mathbb{R})$ is isomorphic to $\mathbb{R}^4$.
The Jacobian of $F$ is
\[ \Jac F(a,b,c,d) = \begin{pmatrix} \partial_a F & \partial_b F & \partial_c F & \partial_d F \end{pmatrix} = \begin{pmatrix} d & -c & -a & b \end{pmatrix}. \]
For $\Jac F$ to not have maximal rank (which is 1), we must have $d = 0$, $c = 0$, $a = 0$, and $b = 0$, which means $A = 0$.
But the zero matrix is excluded from $M$ because its rank is 0, not 1.
Thus $\Jac F(A) \neq 0$ for all $A \in M$. By the regular value theorem, $M$ is a submanifold.
The ambient dimension is $4$ and the number of independent constraints is $1$.
So, the dimension of $M$ is $4 - 1 = 3$.
\end{exercisesolution}

% Source: old_exams/HS18.pdf (Analysis I & II, winter 2018/2019, Prof. Manfred Einsiedler),
% Aufgabe 7, p. 2
% Extractor: Claude Opus 5 (High)
% Kept as a worked example: it is the half of thm:local_parametrization that the theorem's own
% statement hides, namely that the homeomorphism clause is free once one shrinks the domain, and
% that is a piece of theory rather than a drill.
\begin{example}[An immersion is locally a parametrization]
\label{ex:immersion_locally_submanifold}
\Cref{thm:local_parametrization} asks a candidate parametrization $f$ for two things: that $Df_x$
be injective everywhere, and that $f$ be a homeomorphism onto its image. The first is a condition
on a matrix and is easy to check; the second is a global condition and is the one that fails in
practice, as the figure-eight and the lollipop show. The following says that, locally, the second
condition comes for free.

\textbf{Claim.} Let $U \subseteq \mathbb{R}^n$ be open and let $f : U \to \mathbb{R}^{n+1}$ be
smooth with $D_{x_0}f$ injective for every $x_0 \in U$ (such an $f$ is called an
\newterm{immersion} \germanfor{Immersion}). Then every $x_0 \in U$ has an open neighbourhood
$V \subseteq U$ such that $f(V)$ is an $n$-dimensional submanifold of $\mathbb{R}^{n+1}$.

Fix $x_0 \in U$. The matrix $\Jac f(x_0)$ has $n+1$ rows and $n$ columns and is injective, so it
has rank $n$: some $n$ of its rows are linearly independent. Permuting the coordinates of
$\mathbb{R}^{n+1}$ is a linear diffeomorphism of the ambient space, and it carries submanifolds to
submanifolds, so we may assume those are the first $n$ rows. Write accordingly
\[f = \begin{pmatrix} g \\ h \end{pmatrix}, \qquad g : U \to \mathbb{R}^n,\quad
h : U \to \mathbb{R},\]
so that $\Jac g(x_0) \in \mathbb{R}^{n\times n}$ is invertible.

By the inverse function theorem (\cref{thm:inverse_function_theorem}) there is an open
$V \ni x_0$, $V \subseteq U$, such that $g|_V : V \to W := g(V)$ is a diffeomorphism onto an open
set $W \subseteq \mathbb{R}^n$. Let $\varphi := (g|_V)^{-1} : W \to V$ and set
\[\psi := h \circ \varphi : W \to \mathbb{R},\]
which is smooth as a composition of smooth maps. For $x \in V$ we have $g(x) \in W$ and
\[f(x) = \big(g(x),\, h(x)\big) = \big(g(x),\, \psi(g(x))\big),\]
using $\varphi(g(x)) = x$. As $x$ runs over $V$, the point $g(x)$ runs over all of $W$. Therefore
\[f(V) = \{(w,\psi(w)) \mid w \in W\}\]
is exactly the graph of the smooth function $\psi$ over the open set $W$, and a graph is an
$n$-dimensional submanifold of $\mathbb{R}^{n+1}$. To see that last point, apply the regular value
theorem to $(w,s) \mapsto s - \psi(w)$, whose Jacobian $(-\nabla\psi(w)\transp,\ 1)$ never
vanishes.

\begin{remark}[What is local here, and what is not]
The injectivity of $D_{x_0}f$ is a pointwise, linear-algebraic condition, and yet it forces $f$
itself to be injective near $x_0$: on the shrunken $V$ the map $f$ is injective because its first
block $g$ already is. That is the whole trick, and it is why an immersion cannot fail to be a
parametrization for a local reason.

What it cannot repair is a global collision. The figure-eight curve is an immersion of an open
interval whose image meets itself at the crossing point, and no shrinking of the domain changes
that the \emph{image} is not a submanifold there. The claim promises a neighbourhood $V$ in the
\emph{domain} and says nothing about a neighbourhood of $f(x_0)$ in $\mathbb{R}^{n+1}$; the two
are different requests, and only the first is granted. Compare
\cref{ex:lollipop_tangent_field}, where the failure is of exactly this global kind.
\end{remark}
\exinfo{This example is taken from the Analysis I \& II examination of winter 2018/2019
(Prof.\ Manfred Einsiedler), Problem 7, where it is set as a proof exercise. The framing against
\cref{thm:local_parametrization} and the closing remark are added here.}
\end{example}

% Source: old_exams/HS19.pdf (Analysis I & II, 22 January 2020, Prof. Peter S. Jossen),
% Teil B, Aufgabe 4, p. 21
% Extractor: Claude Opus 5 (High)
% Filed next to ex:submanifold_matrices_rank1 deliberately: both live in Mat_2(R) = R^4, and the
% contrast is the point. There the regular value theorem applies as it stands; here the naive
% application fails, because F lands in the symmetric matrices and so is nowhere submersive.
% Parts (e) and (f) of the original invoke the constant rank theorem, which this document does not
% state; the exercise is re-routed through the regular value theorem instead. See the ainote below.
\begin{exercise}[The orthogonal group as a submanifold]
\label{ex:orthogonal_group_submanifold}
Let $\operatorname{Mat}_2(\mathbb{R})$ denote the $2\times 2$ real matrices, identified with
$\mathbb{R}^4$ by listing the entries $a,b,c,d$ of
$\left(\begin{smallmatrix} a & b \\ c & d\end{smallmatrix}\right)$ as a column. Define
\[F : \operatorname{Mat}_2(\mathbb{R}) \to \operatorname{Mat}_2(\mathbb{R}), \qquad
F(A) := AA\transp,\]
and let $\Orth_2(\mathbb{R}) := \{A \in \operatorname{Mat}_2(\mathbb{R}) \mid F(A) = \id\}$.
\begin{enumerate}[label=\textbf{(\alph*)}]
  \item Write $F$ out in the coordinates $a,b,c,d$ and explain why it is smooth.
  \item Show that $F(A)$ is symmetric for every $A$, so that $F$ takes values in the
    three-dimensional subspace $\operatorname{Sym}_2(\mathbb{R}) \subseteq
    \operatorname{Mat}_2(\mathbb{R})$ of symmetric matrices.
  \item Compute the Jacobian matrix of $F$ at $A = \left(\begin{smallmatrix} a & b \\ c &
    d\end{smallmatrix}\right)$. It is a $4\times 4$ matrix.
  \item Show that this Jacobian has rank $3$ at every $A \in \Orth_2(\mathbb{R})$. Why does this
    already rule out applying the regular value theorem to $F$ as it stands?
  \item Regard $F$ instead as a map $\widetilde F : \mathbb{R}^4 \to \mathbb{R}^3$ into
    $\operatorname{Sym}_2(\mathbb{R}) \cong \mathbb{R}^3$ and deduce from
    \cref{thm:regular_value_theorem} that $\Orth_2(\mathbb{R})$ is a submanifold of
    $\operatorname{Mat}_2(\mathbb{R})$. What is its dimension?
  \item Determine the tangent space $T_{\id}\Orth_2(\mathbb{R})$ from the Jacobian computed in
    \textbf{(c)}, and give a basis for it.
\end{enumerate}
\exinfo{This exercise is taken from the Analysis I \& II examination of 22 January 2020
(Prof.\ Peter S.\ Jossen), Part B, Problem 4. The original reaches the submanifold property
through the constant rank theorem, which it also asks the candidate to state; parts \textbf{(d)}
and \textbf{(e)} are rephrased here to use the regular value theorem instead.}
\exsol{sol:orthogonal_group_submanifold}
\end{exercise}

\begin{ainote}
The re-routing in \textbf{(e)} is not a simplification of Jossen's problem, it is the reason the
problem is interesting. The constant rank theorem states that a $C^k$ map whose differential has
the \emph{same} rank $r$ throughout a neighbourhood looks, in suitable charts, like the linear
projection $(x_1,\dots,x_n) \mapsto (x_1,\dots,x_r,0,\dots,0)$; its level sets are then
submanifolds of dimension $n-r$. This document does not state that theorem, and it does not need
to here: restricting the codomain of $F$ to the subspace $\operatorname{Sym}_2(\mathbb{R})$ that
its image already lies in turns the constant rank $3$ into a \emph{full} rank $3$, and the regular
value theorem takes over. That manoeuvre, shrinking the target until the map becomes submersive,
is the standard way to handle matrix groups without the heavier theorem, and it is worth
knowing because $\Orth_n(\mathbb{R})$, $\operatorname{SL}_n(\mathbb{R})$ and their relatives all
fall to it.
\end{ainote}

% Transition: Claude Opus 5 (High)
Three equivalent descriptions are on the table, each characterising the same class of sets. What
they do not yet come with is practice, both in choosing which description to write down for a
surface handed to you and, harder, in arguing that a set is \emph{not} a submanifold at all.
None of the three theorems helps with the second, since each of them only ever yields a positive.
